\begin{slide}
  \begin{block}{Problem}
    Given a multiplication $f_1 \cdot f_2 = p$ of unknown base $b$, what could a possible $b$ be?
  \end{block}
  \pause
  \begin{block}{Insight}
    \begin{itemize}
      \item If $f_1$ and $f_2$ are multiplied in a base greater than $b$, the resulting product
        is always smaller than $p$ in this base.
      \item If $f_1$ and $f_2$ are multiplied in a base smaller than $b$, the resulting product
        is always greater than $p$ in this base.
      \item Since all digits are smaller than or equal to $2^{30}$ and we have at most $1\,000$
        digits, $b$ must be smaller than $2^{61}$.
    \end{itemize}
  \end{block}
\end{slide}
\begin{slide}
  \begin{block}{Solution}
    \begin{itemize} 
      \item Write a multiplication routine that can multiply two numbers in a given base.
      \item Use binary search to determine the correct base.
      \pause
      \item Factorisation approach only leads to accepted solution if a fast factorisation algorithm is used.
    \end{itemize}
  \end{block}
\end{slide}
